Taking an appropriate limit, we see that 
the number of linear extensions of $P \setminus F$ is given by
\begin{equation}
\label{eq:lin_ext}
n! 
\sum_{D \in \EE_P(F)}
 \frac{ 1 }
 { \prod_{v \in P \setminus D} h_P(v) },
\end{equation}
where $n = \# (P \setminus F)$ and $h_P(v)$ is the hook length of $v$ in $P$.
(See Corollary~\ref{cor:main3}\ref{coro5.6_b}.)

If $F = \emptyset$, then our main theorem (Theorem~\ref{thm:main}) gives 
Peterson--Proctor's hook formula (Theorem~\ref{thm:PP}).
If $P = D(\lambda)$ and $F = D(\mu)$ are the Young diagrams of partitions $\lambda \supset \mu$, 
then~\eqref{eq:main} reduces to Morales--Pak--Panova's $q$-hook formula~\cite[Corollary~6.17]{MPP} 
after specializing $z_i = q$ for all $i \in I$, 
and~\eqref{eq:lin_ext} is nothing but Naruse's skew hook formula~\eqref{eq:Naruse-hook}.

Our proof of Theorem~\ref{thm:main} uses the equivariant $K$-theory of Kac--Moody partial flag varieties.
In fact, the notion of excited diagrams and excited peaks for ordinary and shifted Young diagrams 
has its origin in the equivariant Schubert calculus 
(see~\cite{IN1, IN2}; see also~\cite{GK, KN, Kr1, Kr2}).
Theorem~\ref{thm:main} is obtained by proving that 
the both sides of~\eqref{eq:main} equal to the same ratio of the localized equivariant Schubert classes. 
(See Theorems~\ref{thm:xi1} and~\ref{thm:xi2}.)
A key role is played by the Billey-type formula due to Lam--Schilling--Shimozono~\cite{LSS} 
and the Chevalley-type formula due to Lenart--Shimozono~\cite{LS} (see also~\cite{LP}).

This paper is organized as follows.
In Section~\ref{sect2}, we review a definition and basic properties of $d$-complete posets.
In Section~\ref{sect3}, we introduce the notion of excited diagrams for $d$-complete posets, 
which is the key ingredient of the formulation of our main theorem, and study their properties.
In Section~\ref{sect4}, we recall some properties of the equivariant $K$-theory 
and translate the Billey-type formula and the Chevalley-type formula in terms of combinatorics 
of $d$-complete posets.
We will give a proof of our main theorem (Theorem~\ref{thm:main}) and derive some corollaries 
in Section~\ref{sect5}.


\section{%
\texorpdfstring{$d$}{d}-Complete posets
}\label{sect2}

In this section we review a definition and some properties of $d$-complete posets 
and explain their connections to Weyl groups.
See~\cite{Proc1, Proc2, Proc3, Stem2} for details.
We use the terminology for posets (partially ordered sets) 
from~\cite[Chapter~3]{Stan3}.

\subsection{%
Combinatorics of \texorpdfstring{$d$}{d}-complete posets
}

For an integer $k \ge 3$, we denote by $d_k(1)$ the poset consisting of $2k-2$ elements 
$u_1, \ldots, u_{k-2}, x, y, v_{k-2}, \ldots, v_1$ with covering relations
\begin{gather*}
u_1 \gtrdot u_2 \gtrdot \cdots \gtrdot u_{k-2}, \\
u_{k-2} \gtrdot x \gtrdot v_{k-2}, 
\quad u_{k-2} \gtrdot y \gtrdot v_{k-2}, \\
v_{k-2} \gtrdot \cdots \gtrdot v_2 \gtrdot v_1.
\end{gather*}
Note that $x$ and $y$ are incomparable.
The poset $d_k(1)$ is called the \emph{double-tailed diamond}.
The Hasse diagram of $d_k(1)$ is shown in Figure~\ref{fig:diamond}.
\begin{figure}[!htb]
\centering
\setlength{\unitlength}{1.5pt}
\begin{picture}(50,90)(-5,0)
\put(20,5){\circle*{3}}
\put(20,15){\circle*{3}}
\put(20,35){\circle*{3}}
\put(10,45){\circle*{3}}
\put(30,45){\circle*{3}}
\put(20,55){\circle*{3}}
\put(20,75){\circle*{3}}
\put(20,85){\circle*{3}}
\put(20,5){\line(0,1){10}}
\put(20,15){\line(0,1){5}}
\multiput(20,20)(0,2){5}{\line(0,1){1}}
\put(20,30){\line(0,1){5}}
\put(20,35){\line(-1,1){10}}
\put(20,35){\line(1,1){10}}
\put(10,45){\line(1,1){10}}
\put(30,45){\line(-1,1){10}}
\put(20,55){\line(0,1){5}}
\multiput(20,60)(0,2){5}{\line(0,1){1}}
\put(20,70){\line(0,1){5}}
\put(20,75){\line(0,1){10}}
\put(20,80){\makebox(15,10){$u_1$}}
\put(20,70){\makebox(15,10){$u_2$}}
\put(20,50){\makebox(25,10){$u_{k-2}$}}
\put(-5,40){\makebox(15,10){$x$}}
\put(30,40){\makebox(15,10){$y$}}
\put(20,30){\makebox(25,10){$v_{k-2}$}}
\put(20,10){\makebox(15,10){$v_2$}}
\put(20,0){\makebox(15,10){$v_1$}}
\end{picture}
\caption{Double-tailed diamond $d_k(1)$}
\label{fig:diamond}
\end{figure}

Let $P$ be a poset.
An interval $[v,u] = \{ x \in P : v \le x \le u \}$ is called 
a \emph{$d_k$-interval} if it is isomorphic to $d_k(1)$.
Then $v$ and $u$ are called the \emph{bottom} and \emph{top} of $[v,u]$ respectively,
and the two incomparable elements of $[v,u]$ are called the \emph{sides}.
A subset $I$ of $P$ is called \emph{convex} if $x < y < z$ in $P$ and $x$, $z \in I$ imply $y \in I$.
A convex subset $I$ is called a \emph{$d_k^-$-convex set} 
if it is isomorphic to the poset obtained by removing the top element from $d_k(1)$.

\begin{defi}[See~\cite{Okada, PS}]
\label{def:dc_poset}
A poset $P$ is \emph{$d$-complete} if it satisfies 
the following three conditions for every $k \ge 3$:
\begin{enumerate}[(D1)]
\item\label{defi2.1_D1} %[(D1)]
If $I$ is a $d_k^-$-convex set, then there exists an element $u$ 
such that $u$ covers the maximal elements of $I$ and 
$I \cup \{ u \}$ is a $d_k$-interval.
\item\label{defi2.1_D2} %[(D2)]
If $I = [v,u]$ is a $d_k$-interval and the top $u$ covers $u'$ in $P$, 
then $u' \in I$.
\item\label{defi2.1_D3} %[(D3)]
There are no $d_k^-$-convex sets which differ only in the minimal elements.
\end{enumerate}
\end{defi}

It is clear that rooted trees, viewed as posets with their roots being the maximum elements, 
are $d$-complete posets.

\begin{exam}
\label{ex:shape}
For a partition $\lambda$, let $D(\lambda)$ be the Young diagram of $\lambda$ given by
$$
D(\lambda) = \{ (i,j) \in \Int^2 : i \ge 1, \ 1 \le j \le \lambda_i \}.
$$
For a strict partition $\mu$, let $S(\mu)$ be the shifted Young diagram of $\mu$ given by
$$
S(\mu) = \{ (i,j) \in \Int^2 : i \ge 1, \ i \le j \le \mu_i + i-1 \}.
$$
We endow $\Int^2$ with a poset structure by defining 
\begin{equation}
\label{eq:Z2-ordering}
(i,j) \ge (i',j')\text{ if }i \le i'\text{ and }j \le j'.
\end{equation}
Then we regard the Young diagram $D(\lambda)$ and the shifted Young diagram $S(\mu)$ 
as induced subposets of $\Int^2$.
The resulting posets are called a \emph{shape} and a \emph{shifted shape} respectively.
It can be shown that shapes and shifted shapes are $d$-complete posets.
Figure~\ref{fig:Hasse} illustrates the Hasse diagrams of 
$D(5,4,2,1)$ and $S(5,4,2,1)$.
\begin{figure}[!htb]
\centering
%%%%\subcaptionbox{$D(5,4,2,1)$}
{
\setlength{\unitlength}{1.5pt}
\begin{picture}(80,70)(0,-10)
% shape D(5,4,2,1)
\put(5,15){\circle*{3}}
\put(15,25){\circle*{3}}
\put(25,15){\circle*{3}}
\put(25,35){\circle*{3}}
\put(35,25){\circle*{3}}
\put(35,45){\circle*{3}}
\put(45,15){\circle*{3}}
\put(45,35){\circle*{3}}
\put(55,5){\circle*{3}}
\put(55,25){\circle*{3}}
\put(65,15){\circle*{3}}
\put(75,5){\circle*{3}}
\put(5,15){\line(1,1){10}}
\put(15,25){\line(1,-1){10}}
\put(15,25){\line(1,1){10}}
\put(25,15){\line(1,1){10}}
\put(25,35){\line(1,-1){10}}
\put(25,35){\line(1,1){10}}
\put(35,25){\line(1,-1){10}}
\put(35,25){\line(1,1){10}}
\put(35,45){\line(1,-1){10}}
\put(45,15){\line(1,-1){10}}
\put(45,15){\line(1,1){10}}
\put(45,35){\line(1,-1){10}}
\put(55,5){\line(1,1){10}}
\put(55,25){\line(1,-1){10}}
\put(65,15){\line(1,-1){10}}
\put(20,-17){(a) $D(5,4,2,1)$}
\put(110,-17){(b) $S(5,4,2,1)$}
\end{picture}
}
%\subcaptionbox{$S(5,4,2,1)$}
{
\setlength{\unitlength}{1.5pt}
\begin{picture}(80,70)(-15,0)
% shape S(5,4,2,1)
\put(5,5){\circle*{3}}
\put(5,25){\circle*{3}}
\put(5,45){\circle*{3}}
\put(5,65){\circle*{3}}
\put(15,15){\circle*{3}}
\put(15,35){\circle*{3}}
\put(15,55){\circle*{3}}
\put(25,25){\circle*{3}}
\put(25,45){\circle*{3}}
\put(35,15){\circle*{3}}
\put(35,35){\circle*{3}}
\put(45,25){\circle*{3}}
\put(5,5){\line(1,1){10}}
\put(5,25){\line(1,-1){10}}
\put(5,25){\line(1,1){10}}
\put(5,45){\line(1,-1){10}}
\put(5,45){\line(1,1){10}}
\put(5,65){\line(1,-1){10}}
\put(15,15){\line(1,1){10}}
\put(15,35){\line(1,-1){10}}
\put(15,35){\line(1,1){10}}
\put(15,55){\line(1,-1){10}}
\put(25,25){\line(1,-1){10}}
\put(25,25){\line(1,1){10}}
\put(25,45){\line(1,-1){10}}
\put(35,15){\line(1,1){10}}
\put(35,35){\line(1,-1){10}}
\end{picture}
}
\caption{Shape and shifted shape}
\label{fig:Hasse}
\end{figure}
\end{exam}

\begin{exam}
\label{ex:swivel}
Let $P$ be a subset of $\Int^2$ given by
$$
P =
\left\{ 
 \begin{array}{l}
 (1,1), (1,2), (1,3), (1,4), (1,5), (2,3), (2,4), (2,5), \\
 (3,4), (3,5), (3,6), (4,4), (4,5), (4,6), (4,7), (4,8)
 \end{array}
\right\}.
$$
If we regard $P$ as an induced subposet of $\Int^2$ with ordering given by~\eqref{eq:Z2-ordering}, 
then $P$ is a $d$-complete poset, called a \emph{swivel}.
See Figure~\ref{fig:Hasse_swivel} for the Hasse diagram of $P$.
\begin{figure}[!htb]
$$
\setlength{\unitlength}{1.5pt}
%\allinethickness{1.5pt}
\begin{picture}(40,100)
\put(0,100){\circle*{3}}
\put(0,40){\circle*{3}}
\put(10,90){\circle*{3}}
\put(10,70){\circle*{3}}
\put(10,50){\circle*{3}}
\put(10,30){\circle*{3}}
\put(20,80){\circle*{3}}
\put(20,60){\circle*{3}}
\put(20,40){\circle*{3}}
\put(20,20){\circle*{3}}
\put(30,70){\circle*{3}}
\put(30,50){\circle*{3}}
\put(30,30){\circle*{3}}
\put(30,10){\circle*{3}}
\put(40,60){\circle*{3}}
\put(40,0){\circle*{3}}
\put(0,100){\line(1,-1){10}}
\put(0,40){\line(1,1){10}}
\put(0,40){\line(1,-1){10}}
\put(10,90){\line(1,-1){10}}
\put(10,70){\line(1,1){10}}
\put(10,70){\line(1,-1){10}}
\put(10,50){\line(1,1){10}}
\put(10,50){\line(1,-1){10}}
\put(10,30){\line(1,1){10}}
\put(10,30){\line(1,-1){10}}
\put(20,80){\line(1,-1){10}}
\put(20,60){\line(1,1){10}}
\put(20,60){\line(1,-1){10}}
\put(20,40){\line(1,1){10}}
\put(20,40){\line(1,-1){10}}
\put(20,20){\line(1,1){10}}
\put(20,20){\line(1,-1){10}}
\put(30,70){\line(1,-1){10}}
\put(30,50){\line(1,1){10}}
\put(30,10){\line(1,-1){10}}
\end{picture}
$$
\caption{Swivel}
\label{fig:Hasse_swivel}
\end{figure}
\end{exam}

A poset $P$ is called \emph{connected} if its Hasse diagram is a connected graph.
It is easy to see that, if $P$ is a $d$-complete poset, 
then each connected component of $P$ is $d$-complete.
(For our purpose to prove Theorem~\ref{thm:main}, 
there is no harm in assuming that a $d$-complete poset is connected.)

\newbox\toto
\setbox\toto=\hbox{\cite[\S~3]{Proc1}}
\begin{prop}[\box\toto]
\label{prop:dc_connected}
%\textup{(\cite[\S~3]{Proc1})}
Let $P$ be a $d$-complete poset.
If $P$ is connected, then $P$ has a unique maximal element.
\end{prop}

Let $P$ be a finite poset.
The \emph{top forest} $\Gamma$ of $P$ is a full subgraph of 
the Hasse diagram of $P$, whose vertex set consists of all elements $x \in P$ 
such that every $y \ge x$ is covered by at most one other element. 
Note that two elements $x$ and $y \in \Gamma$ are connected by an edge in $\Gamma$ 
if and only if $x$ covers $y$ or $y$ covers $x$.
Then $\Gamma$ becomes a forest as a graph (i.e. a disjoint union of trees).
If $P$ is a connected $d$-complete poset, then $\Gamma$ is a tree called 
the \emph{top tree} of $P$.
We can regard the top forest $\Gamma$ as a Dynkin diagram of a Kac--Moody group.
For example, the top trees of Figure~\ref{fig:Hasse} (a) and (b) are
of type $A_8$ and $D_6$ respectively. 
If $P$ is the swivel in Figure~\ref{fig:Hasse_swivel}, the top tree is
the Dynkin diagram of type $E_6$.
(See also Subections~\ref{subsec2.2} and~\ref{subsec4.3}.)

Remark~\ref{rem:heap1} below will indicate how to combine the results 
of~\cite{Proc2} and~\cite{Stem2} to obtain the following statement.

\setbox\toto=\hbox{\cite[Proposition~8.6]{Proc2},~\cite[Proposition~3.1]{Stem2}}
\begin{prop}[\box\toto]
\label{prop:dc_color}
%\textup{(\cite[Proposition~8.6]{Proc2},~\cite[Proposition~3.1]{Stem2})}
Let $P$ be a $d$-complete poset and $\Gamma$ its top forest.
Let $I$ be a set of colors whose cardinality is the same as $\Gamma$.
Then a bijective labeling $c : \Gamma \to I$ can be uniquely extended to a map 
$c : P \to I$ satisfying the following three conditions:
\begin{enumerate}[(C1)]
\item\label{prop2.5_C1} %[(C1)]
If $x$ and $y$ are incomparable, then $c(x) \neq c(y)$.
\item\label{prop2.5_C2} %[(C2)]
If an interval $[v,u]$ is a chain, then the colors $c(x)$ $(x \in [v,u])$ are 
distinct.
\item\label{prop2.5_C3} %[(C3)]
If $[v,u]$ is a $d_k$-interval then $c(v) = c(u)$.
\end{enumerate}
Moreover this map $c$ satisfies
\begin{enumerate}[(C1)]\setcounter{enumi}{3}
\item\label{prop2.5_C4} %[(C4)]
If $x$ covers $y$, then the nodes labeled by $c(x)$ and $c(y)$ are adjacent in $\Gamma$.
\item\label{prop2.5_C5} %[(C5)]
If $c(x) = c(y)$ or the nodes labeled by $c(x)$ and $c(y)$ are adjacent in $\Gamma$, 
then $x$ and $y$ are comparable.
\end{enumerate}
Such a map $c : P \to I$ is called a \emph{$d$-complete coloring}.
\end{prop}

Let $P$ be a $d$-complete poset and 
$c : P \to I$ a $d$-complete coloring.
Let $\vectz = (z_i)_{i \in I}$ be indeterminates.
Given an order filter $F$ of $P$, we regard $P \setminus F$ as an induced subposet.
For a $(P \setminus F)$-partition $\sigma \in \AP(P \setminus F)$, we put
$$
\vectz^\sigma = \prod_{v \in P \setminus F} z_{c(v)}^{\sigma(v)}.
$$
We are interested in the multivariate generating function
$$
\sum_{\sigma \in \AP(P \setminus F)} \vectz^\sigma
$$
of $(P \setminus F)$-partitions.
For a subset $D$ of $P$, we write
$$
\vectz[D] = \prod_{v \in D} z_{c(v)}.
$$
Instead of giving a definition of hooks $H_P(u) \subset P$ for a general 
$d$-complete poset $P$, we define associated monomials $\vectz[H_P(u)]$ 
directly by induction as follows:

\begin{defi}
\label{def:hook}
Let $P$ be a $d$-complete poset with $d$-complete coloring $c : P \to I$.
\begin{enumerate}[(i)]
\item\label{defi2.6_i} %[(i)]
If $u$ is not the top of any $d_k$-interval, then we define
$$
\vectz[H_P(u)] = \prod_{w \le u} z_{c(w)}.
$$
\item\label{defi2.6_ii} %[(ii)]
If $u$ is the top of a $d_k$-interval $[v,u]$, then we define
$$
\vectz[H_P(u)]
 =
\frac{ \vectz[H_P(x)] \cdot \vectz[H_P(y)] }
 { \vectz[H_P(v)] },
$$
where $x$ and $y$ are the sides of $[v,u]$.
\end{enumerate}
\end{defi}

\begin{exam}
\label{ex:shape_hc}
Let $P = D(\lambda)$ be the shape corresponding to a partition $\lambda$.
Then the top tree $\Gamma$ of $D(\lambda)$ is given by
$$
\Gamma = \{ (1,j) : 1 \le j \le \lambda_1 \} \cup \{ (i,1) : 1 \le i \le \lambda'_1 \},
$$
where $\lambda'_1$ is the number of cells in the first column of the Young diagram $D(\lambda)$.
A $d$-complete coloring $c: D(\lambda) \to I = \{ -(\lambda'_1-1), \dots, -1, 0, 1, \dots, \lambda_1-1 \}$ 
is given by
$$
c(i,j) = j-i.
$$
The classical definition of the hook $H_{D(\lambda)}(i,j) \subset D(\lambda)$ 
at $u = (i,j)$ in $D(\lambda)$ is as follows:
$$
H_{D(\lambda)}(i,j)
 =
\{ (i,j) \}
\cup
\{ (i,l) \in D(\lambda) : l > j \}
\cup
\{ (k,j) \in D(\lambda) : k > i \}.
$$
Then the hook monomial $\vectz[H_{D(\lambda)}(i,j)]$ in Definition~\ref{def:hook} 
and the hook $H_{D(\lambda)}(i,j)$ are related as 
$$
\vectz[H_{D(\lambda)}(i,j)]
 =
\prod_{(p,q)\in H_{D(\lambda)}(i,j) } z_{c(p,q)}.
$$
\end{exam}

\begin{exam}
\label{ex:shifted_hc}
Let $P = S(\mu)$ be the shifted shape corresponding to a strict partition $\mu$ of length $\ge 2$.
Then the top tree $\Gamma$ of $S(\mu)$ is given by
$$
\Gamma = \{ (1,j) : 1 \le j \le \mu_1 \} \cup \{ (2, 2) \},
$$
and a $d$-complete coloring $c : S(\mu) \to I = \{ 0, 0', 1, 2, \dots, \mu_1 - 1 \}$ is given by
$$
c(i,j)
 =
\begin{cases}
 j-i &\text{if $i < j$,} \\
 0 &\text{if $i=j$ and $i$ is odd,} \\
 0' &\text{if $i=j$ and $i$ is even.}
\end{cases}
$$
The (shifted) hook $H_{S(\mu)}(i,j)$ of $u = (i,j)$ in $S(\mu)$ is the subset of $S(\mu)$ defined by
\begin{align*}
H_{S(\mu)}(i,j)
 &=
\{ (i,j) \}
\cup
\{ (i,l) \in S(\mu) : l > j \}
\cup
\{ (k,j) \in S(\mu) : k > j \}
\\
&\quad
\cup
\{ (j+1,l) \in S(\mu) : l > j \}.
\end{align*}
For example, if $\mu = (5,4,2,1)$ and $(i,j) = (1,2)$, then the corresponding hook is given by
$$
H_{S(5,4,2,1)}(1,2)
 =
\{ (1,2), (1,3), (1,4), (1,5), (2,2), (3,3), (3,4) \},
$$
and the hook monomial is 
$\vectz[H_{S(5,4,2,1)}(1,2)]= z_{0'}z_{0} z_1^2 z_2 z_3 z_4$.
\end{exam}

\subsection{%
\texorpdfstring{$d$}{d}-Complete posets and Weyl groups
}\label{subsec2.2}

Let $P$ be a connected $d$-complete poset with top tree $\Gamma$.
We regard $\Gamma$ as a (simply-laced) Dynkin diagram with node set $I$ 
and the $d$-complete coloring as a map $c : P \to I$.
Let $A = (a_{ij})_{i,j \in I}$ be the generalized Cartan matrix of $\Gamma$ given by
$$
a_{ij}
 =
\begin{cases}
 2 &\text{if $i=j$,} \\
 -1 &\text{if $i \neq j$ and $i$ and $j$ are adjacent in $\Gamma$,} \\
 0 &\text{otherwise.}
\end{cases}
$$
We fix the following data associated to $A$:
\begin{itemize}
\item
a free $\Int$-module $\Lambda$ of finite rank at least $\# \Gamma$, called the \emph{weight lattice},
\item
a linearly independent subset $\Pi = \{ \alpha_i : i \in I \}$ of $\Lambda$, called the \emph{simple roots},
\item
a subset $\Pi^\vee = \{ \alpha^\vee_i : i \in I \}$ of the dual lattice $\Lambda^* = \Hom_\Int(\Lambda,\Int)$, 
called the \emph{simple coroots},
\item
a subset $\{ \lambda_i : i \in I \}$ of $\Lambda$, called the \emph{fundamental weights},
\end{itemize}
such that
$$
\langle \alpha^\vee_i, \alpha_j \rangle = a_{ij},
\quad
\langle \alpha^\vee_i, \lambda_j \rangle = \delta_{ij},
$$
where $\langle \quad, \quad \rangle : \Lambda^* \times \Lambda \to \Int$ is the canonical pairing.
Let $W$ be the corresponding 
(Kac--Moody) Weyl group generated by the simple reflections $\{ s_i : i \in I \}$, 
where $s_i$ acts on $\Lambda$ and $\Lambda^*$ by the rule
$$
s_i(\lambda) = \lambda - \langle \alpha_i^\vee, \lambda \rangle \alpha_i
\quad(\lambda \in \Lambda),
\quad
s_i(\lambda^\vee) = \lambda^\vee - \langle \lambda^\vee, \alpha_i \rangle \alpha^\vee_i
\quad(\lambda^\vee \in \Lambda^*).
$$
Then $W$ is a Coxeter group, and we have the length function $l$ and the Bruhat order $<$ on $W$.
The set of real roots $\Phi$ and the set of real coroots $\Phi^\vee$ are defined by 
$\Phi = W \Pi$ and $\Phi^\vee = W \Pi^\vee$ respectively.
The set of simple roots $\Pi$ (\resp the set of simple coroots $\Pi^\vee$) determines 
the decomposition of $\Phi$ (\resp $\Phi^\vee$) into the positive system $\Phi_+$ (\resp $\Phi^\vee_+$) and 
the negative system $\Phi_-$ (\resp $\Phi^\vee_-$).
We introduce the standard partial ordering on $\Phi_+$ (\resp $\Phi^\vee_+$) by setting 
$\alpha > \beta$ if $\alpha - \beta$ is a sum of simple roots $\{ \alpha_i : i \in I \}$ 
(\resp $\alpha^\vee > \beta^\vee$ if $\alpha^\vee - \beta^\vee$ is a sum of simple coroots 
$\{ \alpha^\vee_i : i \in I \}$).

For $p \in P$, we put
$$
\alpha(p) = \alpha_{c(p)},
\quad
\alpha^\vee(p) = \alpha^\vee_{c(p)},
\quad
s(p) = s_{c(p)}.
$$
Let $\alpha_P$ and $\lambda_P$ be the simple root and the fundamental weight 
corresponding to the color $i_P$ of the maximum element of $P$. 

Take a linear extension and label the elements of $P$ with $p_1, \ldots, p_N$ ($N= \# P$) 
so that $p_i < p_j$ in $P$ implies $i<j$.
Then we construct an element $w_P \in W$ by putting
$$
w_P = s(p_1) s(p_2) \cdots s(p_N).
$$
A Weyl group element $w \in W$ is called \emph{$\lambda$-minuscule} 
if there exists a reduced expression $w = s_{i_1} \cdots s_{i_l}$ such that
$$
\langle \alpha^\vee_{i_k}, s_{i_{k+1}} \cdots s_{i_l} \lambda \rangle = 1
\quad(1 \le k \le l),
$$
or equivalently
$$
s_{i_k} \cdots s_{i_l} \lambda = \lambda - \alpha_{i_k} - \cdots - \alpha_{i_l}.
$$
A element $w \in W$ is called \emph{fully commutative} 
if any reduced expression of $w$ can be obtained from any other by using only 
the Coxeter relations of the form $st = ts$.

\setbox\toto=\hbox{See~\cite{Proc2} and~\cite[Proposition~2.1]{Stem2}}
\begin{prop}[\box\toto]
\label{prop:w_P}
%\textup{(See~\cite{Proc2} and~\cite[Proposition~2.1]{Stem2})}
Let $P$ be a connected $d$-com\-plete poset.
Then the Weyl group element $w_P \in W$ is $\lambda_P$-minuscule and hence fully commutative.
\end{prop}

If $p = p_k \in P$, then we define 
\begin{align*}
\beta(p_k) &= s(p_1) \cdots s(p_{k-1}) \alpha(p_k),
\\
\gamma(p_k) &= s(p_N) \cdots s(p_{k+1}) \alpha(p_k),
\\
\gamma^\vee(p_k) &= s(p_N) \cdots s(p_{k+1}) \alpha^\vee(p_k).
\end{align*}
It follows from Proposition~\ref{prop:w_P} that, 
for each $p \in P$, the roots $\beta(p)$, $\gamma(p)$ and the coroot $\gamma^\vee(p)$ 
are independent of the choices of linear extensions.
For a Weyl group element $w \in W$, we put
$$
\Phi(w) = \Phi_+ \cap w \Phi_-,
\quad
\Phi^\vee(w) = \Phi^\vee_+ \cap w \Phi^\vee_-.
$$
Then it is well-known (see~\cite[\S~5.6]{H}) that
$$
\Phi(w_P) = \{ \beta(p) : p \in P \},
\quad
\Phi^\vee(w_P^{-1}) = \{ \gamma^\vee(p) : p \in P \}.
$$
Moreover we have

\begin{prop}
\label{prop:root}
Let $P$ be a connected $d$-complete poset. Then we have
\begin{enumerate}[(a)]
\item\label{prop2.10_a} %[(a)]
\textup{(See~\cite[Proposition~3.1 and Theorem~5.5]{Stem2})}
The poset $P$ is isomorphic to the order dual of $\Phi^\vee(w_P^{-1})$ 
with the standard coroot ordering on $\Phi^\vee_+$.
\item\label{prop2.10_b} %[(b)]
\textup{(See~\cite[Lemma~IV]{Proc3})}
Under the identification $z_i = e^{\alpha_i}$ \textup{(}$i \in I$\textup{)}, we have
$$
\vectz[H_P(p)] = e^{\beta(p)}
\quad(p \in P).
$$
\item\label{prop2.10_c} %[(c)]
\textup{(See~\cite[Proposition~5.1]{Stem2})}
We have
$$
\langle \gamma^\vee(p), \lambda_P \rangle = 1
\quad(p \in P).
$$
\end{enumerate}
\end{prop}

Let $W_{\lambda_P}$ be the stabilizer of $\lambda_P$ in $W$.
Then $W_{\lambda_P}$ is the maximal parabolic subgroup corresponding to $I \setminus \{ i_P \}$.
Let $W^{\lambda_P}$ be the set of minimum length coset representatives of $W/W_{\lambda_P}$.
For a subset $D = \{ p_{i_1}, \ldots, p_{i_r} \}$ ($i_1 < \cdots < i_r$) of $P$, we define
\begin{equation}
\label{eq:w_D}
w_D = s(p_{i_1}) \cdots s(p_{i_r}).
\end{equation}
Since $w_P$ is fully commutative (Proposition~\ref{prop:w_P}), 
we see that $w_D$ is independent of the choices of linear extensions of $P$.

\begin{prop}
\label{prop:Weyl}
Let $P$ be a connected $d$-complete poset.
Then we have
\begin{enumerate}[(a)]
\item\label{prop2.11_a} %[(a)]
\textup{(See~\cite[Proposition~I]{Proc3})}
The map $F \mapsto w_F$ gives a poset isomorphism 
from the set of all order filters of $P$ ordered by inclusion 
to the Bruhat interval $[e,w_P]$ in $W^{\lambda_P}$.
\item\label{prop2.11_b} %[(b)]
\textup{(See~\cite[Remark~2.7~(b)]{Stem2})}
If $F$ is an order filter of $P$, then $w_F$ is $\lambda_P$-minuscule, and
$$
w_F \lambda_P = \lambda_P - \sum_{p \in F} \alpha(p).
$$
\end{enumerate}
\end{prop}

\begin{rema}
\label{rem:heap1}
Let $A = (a_{ij})_{i,j \in I}$ be a symmetrizable generalized Cartan matrix, 
and $\Gamma$ the corresponding Dynkin diagram with node set $I$.
Let $W$ be the associated Kac--Moody Weyl group generated by $\{ s_i : i \in I \}$.
Given a (not necessarily reduced) expression $s_{i_1} s_{i_2} \dots s_{i_N}$ of an element $w \in W$ 
in simple reflections, we can define a poset $H$, called the \emph{heap}, as follows 
(see~\cite{Stem1}).
The poset $H$ consists of the ground set $\{ 1, 2, \dots, N \}$ and the partial ordering 
$\preceq$ obtained by taking the transitive closure of the relations given by
$$
\text{$a \prec b$ if $a < b$ and either $s_{i_a} s_{i_b} \neq s_{i_b} s_{i_a}$ or $i_a = i_b$.}
$$
The heap $H$ has a natural labeling (coloring) $c: H \to I$ given by $c(a) = i_a$.
If $w \in W$ is fully commutative, then the heap defined by a reduced expression of $w$ 
is independent of the choices of reduced expressions.
In this case we denote the resulting heap by $H(w)$.

If $a_{ij} = a_{ji} \in \{ 0, -1 \}$ for any pair of vertices $i$, $j \in I$ ($i \neq j$), 
then we say that $A$, $W$ and $H(w)$ are \emph{simply-laced}; 
otherwise they are \emph{multiply-laced}.
For a dominant weight $\lambda$ and a $\lambda$-minuscule element $w$ in the simply-laced Weyl group $W$,
Proctor~\cite{Proc2} proved that the interval $[e,w]$ of $W^\lambda$ with respect to 
the Bruhat order is a distributive lattice, 
and that the order dual of the induced subposet consisting of all join-irreducible elements of $[e,w]$ 
is a $d$-complete poset colored by $I$.
Stembridge~\cite{Stem2} extends Proctor's result to any symmetrizable Kac--Moody Weyl group 
in the setting of heaps.
The results of~\cite{Proc2} and~\cite{Stem2} may be combined by using 
\cite[Theorem~3.2~(c)]{Stem1} to produce Proposition~\ref{prop:dc_color} above.
Henceforth the simply-laced case will be referred to with the ``$d$-complete poset'' terminology 
and the general case will be referred to with the ``heap'' terminology.
Every $d$-complete poset $P$ is isomorphic to the simply-laced heap $H(w_P)$. 
In general, if $w \in W$ is dominant minuscule, i.e. $\lambda$-minuscule for some dominant weight $\lambda$, 
the corresponding heap $H(w)$ is isomorphic (as a unlabeled poset) to a $d$-complete poset.
See~\cite[Sections~3 and~4]{Stem2}.

The propositions in this section hold literally for heaps $H(w)$ of dominant minuscule elements, 
except for Proposition~\ref{prop:dc_color} and Proposition~\ref{prop:root}~\ref{prop2.10_b}.
The latter half of Proposition~\ref{prop:dc_color} holds for heaps, 
i.e. the labeling $c : H(w) \to I$ satisfies~\ref{prop2.5_C4} and~\ref{prop2.5_C5}.
And we adopt Proposition~\ref{prop:root}~\ref{prop2.10_b} as a definition of the hook monomial for $H(w)$.
Below is an example of a non-simply-laced heap.
\end{rema}

\begin{exam} (Non-simply-laced heap)
\label{ex:shifted-B}
Let $\mu = (\mu_1, \dots, \mu_l)$ be a strict partition of length $l$.
Then the shifted shape $S(\mu)$ can be regarded as 
the heap associated to a dominant minuscule element of the Weyl group of type $B$,
which is not simply-laced.
Put $m = \mu_1$ and let $W$ be the Weyl group generated by $s_0, s_1, \dots, s_{m-1}$ 
subject to the relations
\begin{gather*}
s_i^2 = 1 \quad(i=0, 1, \dots, m-1),
\\
s_0 s_1 s_0 s_1 = s_1 s_0 s_1 s_0,
\\
s_i s_{i+1} s_i = s_{i+1} s_i s_{i+1} \quad(i=1, 2, \dots, m-2),
\\
s_i s_j = s_j s_i \quad(|i-j| \ge 2).
\end{gather*}
Then we define an element $w_\mu \in W$ by putting
$$
w_\mu = (s_{\mu_l-1} \cdots s_1 s_0) \cdots (s_{\mu_2-1} \cdots s_1 s_0) (s_{\mu_1-1} \cdots s_1 s_0).
$$
Then it can be shown that $w_\mu$ is $\lambda_0$-minuscule, where $\lambda_0$ is the fundamental weight 
corresponding to $s_0$, and that the map
\begin{equation}
\label{eq:S=H}
S(\mu) \ni (i,j) \mapsto \mu_l + \dots + \mu_i - j + i \in H(w_\mu)
\end{equation}
gives a poset isomorphism.
We identify the ground set of $H(w_\mu)$ with $S(\mu)$ via the isomorphism~\eqref{eq:S=H}.
Then the natural labeling $c' : S(\mu) \to \{ 0, 1, \dots, m-1 \}$ of $H(w_\mu)$ is given by 
$c'(i,j) = j-i$, which is different from the $d$-complete coloring of $S(\mu)$ given 
in Example~\ref{ex:shifted_hc}.
And the ``hook'' $H'_{S(\mu)}(v)$ (see~\cite[Definition~4.8]{N2}) is defined by
\begin{align*}
H'_{S(\mu)}(i,j)
 &=
\{ (i,j) \}
 \cup
\{ (i,l) \in S(\mu) : l > j \}
 \cup
\{ (k,j) \in S(\mu) : k > i \}
\\
 &\quad
\cup
\begin{cases}
 \{ (i,i) \} \cup \{ (j,l) \in S(\mu) : l \ge j \} &\text{if $i<j$ and $(j,j) \in S(\mu)$,} \\
 \emptyset &\text{otherwise,}
\end{cases}
\end{align*}
which is also different from the shifted hook given in Example~\ref{ex:shifted_hc}.
(See also Example~\ref{ex:heap}.)
\end{exam}

\section{%
Excited diagrams
}\label{sect3}

In this section we introduce the notion of excited and $K$-theoretical excited diagrams 
in a $d$-complete poset and study their properties.

\subsection{%
Excited diagrams
}

First we generalize the notion of exited diagrams for Young diagram and shifted Young diagram, 
which were introduced by Ikeda--Naruse~\cite{IN1} and Kreiman~\cite{Kr1, Kr2} independently, 
to a general $d$-complete posets.
And we give a generalization of backward movable positions or excited peaks 
introduced in~\cite{IN2, KN, MPP}.

Let $P$ be a $d$-complete poset with top forest $\Gamma$ and $d$-complete coloring $c : P \to I$.
For a subset $D \subset P$ and a color $i \in I$, we put
$$
D_i = \{ x \in D : c(x) = i \}.
$$
For $i \in I$, let $N_i$ be the subset of $P$ consisting of element $x \in P$ whose color 
$c(x)$ is adjacent to $i$ in the Dynkin diagram $\Gamma$.
Note that, if $[v,u]$ is a $d_k$-interval, then $[v,u] \cap N_{c(u)}$ consists of elements $x \in [v,u]$ 
such that $x$ is covered by $u$ or covers $v$.

\begin{defi}
\label{def:excited}
Let $P$ be a $d$-complete poset and let $F$ be an order filter of $P$.
\begin{enumerate}[(a)]
\item\label{defi3.1_a} %[(a)]
Let $D$ be a subset of $P$ and $u \in D$.
We say that $u$ is \emph{$D$-active} 
if there exists an element $v \in (P \setminus D)_{c(u)}$ 
such that $v < u$, $[v,u]$ is a $d_k$-interval and
$$
[v,u] \cap D \cap N_{c(u)} = \emptyset.
$$
\item\label{defi3.1_b} %[(b)]
Let $D$ be a subset of $P$ and $u \in D$.
If $u$ is $D$-active, then we define $\alpha_u(D)$ to be the subset of $P$ obtained from $D$ by 
replacing $u \in D$ by the bottom element $v$ of the $d_k$-interval $[v,u]$.
We call this replacement an (ordinary) \emph{elementary excitation}.
\item\label{defi3.1_c} %[(c)]
An \emph{excited diagram} of $F$ in $P$ is a subset of $P$ 
obtained from $F$ after a sequence of elementary excitations on active elements. 
Let $\EE_P(F)$ be the set of all excited diagrams of $F$ in $P$.
\item\label{defi3.1_d} %[(d)]
To an excited diagram $D \in \EE_P(F)$ we associate a subset $B(D) \subset P$ as follows:
If $D = F$, then $B(F) = \emptyset$.
If $D$ is an excited diagram with an active element $u$, then we define
$$
B(\alpha_u(D))
 = 
\left( B(D) \setminus ([v,u] \cap N_{c(u)}) \right) \cup \{ u \},
$$
where $[v,u]$ is the $d_k$-interval with top element $u$.
We call $B(D)$ the set of \emph{excited peaks} of $D$.
(We will show that $B(D)$ is a well-defined subset of $P \setminus D$ in Proposition~\ref{prop:BD}.)
\end{enumerate}
\end{defi}

In general, if two elements $u \in D$ and $v \not\in D$ with $v < u$ 
have the same color $i = c(v) = c(u)$ and satisfy $[v,u] \cap D \cap N_i = \emptyset$, 
then $D \setminus \{ u \} \cup \{ v \}$ is obtained from $D$ by a sequence 
of elementary excitations.

If $P$ is a shape or a shifted shape, our definition above coincides with the definitions of 
elementary excitations in~\cite{IN1, IN2}, 
or ladder moves in~\cite{Kr1, Kr2}, 
and backward movable positions in~\cite{IN2, KN} 
or excited peaks in~\cite{MPP} (only for a shape).

\begin{exam}
\label{ex:shape_excited}
If $P = D(5,4,2,1)$ is the shape corresponding to a partition $(5,4,2,1)$ and 
$F = D(3,1)$, then there are $7$ excited diagrams in $\EE_P(F)$ shown in Figure~\ref{fig:shape_excited}.
In Figure~\ref{fig:shape_excited} (and Figures~\ref{fig:swivel_excited}) 
the shaded cells form an exited diagram and a cell with $\times$ is an excited peak.
And the arrow $D \longrightarrow D'$ means that $D'$ is obtained from $D$ by an elementary excitation.
\begin{figure}[t]
$$
\setlength{\unitlength}{1.1pt}
\begin{CD}
\raisebox{-20pt}{
\begin{picture}(50,40)
\fboxsep=0mm
\put(0,30){\colorbox[gray]{0.7}{\makebox(10,10){}}}
\put(10,30){\colorbox[gray]{0.7}{\makebox(10,10){}}}
\put(20,30){\colorbox[gray]{0.7}{\makebox(10,10){}}}
\put(0,20){\colorbox[gray]{0.7}{\makebox(10,10){}}}
\put(0,40){\line(1,0){50}}
\put(0,30){\line(1,0){50}}
\put(0,20){\line(1,0){40}}
\put(0,10){\line(1,0){20}}
\put(0,0){\line(1,0){10}}
\put(0,0){\line(0,1){40}}
\put(10,0){\line(0,1){40}}
\put(20,10){\line(0,1){30}}
\put(30,20){\line(0,1){20}}
\put(40,20){\line(0,1){20}}
\put(50,30){\line(0,1){10}}
\end{picture}
}
@>>>
\raisebox{-20pt}{
\begin{picture}(50,40)
\fboxsep=0mm
\put(0,30){\colorbox[gray]{0.7}{\makebox(10,10){}}}
\put(10,30){\colorbox[gray]{0.7}{\makebox(10,10){}}}
\put(30,20){\colorbox[gray]{0.7}{\makebox(10,10){}}}
\put(0,20){\colorbox[gray]{0.7}{\makebox(10,10){}}}
\put(0,40){\line(1,0){50}}
\put(0,30){\line(1,0){50}}
\put(0,20){\line(1,0){40}}
\put(0,10){\line(1,0){20}}
\put(0,0){\line(1,0){10}}
\put(0,0){\line(0,1){40}}
\put(10,0){\line(0,1){40}}
\put(20,10){\line(0,1){30}}
\put(30,20){\line(0,1){20}}
\put(40,20){\line(0,1){20}}
\put(50,30){\line(0,1){10}}
\put(20,30){\makebox(10,10){$\times$}}
\end{picture}
}
@>>>
\raisebox{-20pt}{
\begin{picture}(50,40)
\fboxsep=0mm
\put(0,30){\colorbox[gray]{0.7}{\makebox(10,10){}}}
\put(20,20){\colorbox[gray]{0.7}{\makebox(10,10){}}}
\put(30,20){\colorbox[gray]{0.7}{\makebox(10,10){}}}
\put(0,20){\colorbox[gray]{0.7}{\makebox(10,10){}}}
\put(0,40){\line(1,0){50}}
\put(0,30){\line(1,0){50}}
\put(0,20){\line(1,0){40}}
\put(0,10){\line(1,0){20}}
\put(0,0){\line(1,0){10}}
\put(0,0){\line(0,1){40}}
\put(10,0){\line(0,1){40}}
\put(20,10){\line(0,1){30}}
\put(30,20){\line(0,1){20}}
\put(40,20){\line(0,1){20}}
\put(50,30){\line(0,1){10}}
\put(10,30){\makebox(10,10){$\times$}}
\end{picture}
}
\\[5pt]
@VVV
@VVV
@VVV
\\[5pt]
\raisebox{-20pt}{
\begin{picture}(50,40)
\fboxsep=0mm
\put(0,30){\colorbox[gray]{0.7}{\makebox(10,10){}}}
\put(10,30){\colorbox[gray]{0.7}{\makebox(10,10){}}}
\put(20,30){\colorbox[gray]{0.7}{\makebox(10,10){}}}
\put(10,10){\colorbox[gray]{0.7}{\makebox(10,10){}}}
\put(0,40){\line(1,0){50}}
\put(0,30){\line(1,0){50}}
\put(0,20){\line(1,0){40}}
\put(0,10){\line(1,0){20}}
\put(0,0){\line(1,0){10}}
\put(0,0){\line(0,1){40}}
\put(10,0){\line(0,1){40}}
\put(20,10){\line(0,1){30}}
\put(30,20){\line(0,1){20}}
\put(40,20){\line(0,1){20}}
\put(50,30){\line(0,1){10}}
\put(0,20){\makebox(10,10){$\times$}}
\end{picture}
}
@>>>
\raisebox{-20pt}{
\begin{picture}(50,40)
\fboxsep=0mm
\put(0,30){\colorbox[gray]{0.7}{\makebox(10,10){}}}
\put(10,30){\colorbox[gray]{0.7}{\makebox(10,10){}}}
\put(30,20){\colorbox[gray]{0.7}{\makebox(10,10){}}}
\put(10,10){\colorbox[gray]{0.7}{\makebox(10,10){}}}
\put(0,40){\line(1,0){50}}
\put(0,30){\line(1,0){50}}
\put(0,20){\line(1,0){40}}
\put(0,10){\line(1,0){20}}
\put(0,0){\line(1,0){10}}
\put(0,0){\line(0,1){40}}
\put(10,0){\line(0,1){40}}
\put(20,10){\line(0,1){30}}
\put(30,20){\line(0,1){20}}
\put(40,20){\line(0,1){20}}
\put(50,30){\line(0,1){10}}
\put(0,20){\makebox(10,10){$\times$}}
\put(20,30){\makebox(10,10){$\times$}}
\end{picture}
}
@>>>
\raisebox{-20pt}{
\begin{picture}(50,40)
\fboxsep=0mm
\put(0,30){\colorbox[gray]{0.7}{\makebox(10,10){}}}
\put(20,20){\colorbox[gray]{0.7}{\makebox(10,10){}}}
\put(30,20){\colorbox[gray]{0.7}{\makebox(10,10){}}}
\put(10,10){\colorbox[gray]{0.7}{\makebox(10,10){}}}
\put(0,40){\line(1,0){50}}
\put(0,30){\line(1,0){50}}
\put(0,20){\line(1,0){40}}
\put(0,10){\line(1,0){20}}
\put(0,0){\line(1,0){10}}
\put(0,0){\line(0,1){40}}
\put(10,0){\line(0,1){40}}
\put(20,10){\line(0,1){30}}
\put(30,20){\line(0,1){20}}
\put(40,20){\line(0,1){20}}
\put(50,30){\line(0,1){10}}
\put(0,20){\makebox(10,10){$\times$}}
\put(10,30){\makebox(10,10){$\times$}}
\end{picture}
}
@>>>
\raisebox{-20pt}{
\begin{picture}(50,40)
\fboxsep=0mm
\put(10,20){\colorbox[gray]{0.7}{\makebox(10,10){}}}
\put(20,20){\colorbox[gray]{0.7}{\makebox(10,10){}}}
\put(30,20){\colorbox[gray]{0.7}{\makebox(10,10){}}}
\put(10,10){\colorbox[gray]{0.7}{\makebox(10,10){}}}
\put(0,40){\line(1,0){50}}
\put(0,30){\line(1,0){50}}
\put(0,20){\line(1,0){40}}
\put(0,10){\line(1,0){20}}
\put(0,0){\line(1,0){10}}
\put(0,0){\line(0,1){40}}
\put(10,0){\line(0,1){40}}
\put(20,10){\line(0,1){30}}
\put(30,20){\line(0,1){20}}
\put(40,20){\line(0,1){20}}
\put(50,30){\line(0,1){10}}
\put(0,30){\makebox(10,10){$\times$}}
\end{picture}
}
\end{CD}
$$
\caption{Excited diagrams of $D(3,1)$ in $D(5,4,2,1)$}
\label{fig:shape_excited}
\end{figure}
\end{exam}

\begin{exam}
\label{ex:swivel_excited}
If $P$ is the swivel given in Example~\ref{ex:swivel} and $F$ is the order filter consisting of three elements, 
then there are $8$ excited diagrams in $\EE_P(F)$ shown in Figure~\ref{fig:swivel_excited}.
\begin{figure}[!htb]
$$
\setlength{\unitlength}{1.1pt}
\begin{CD}
\raisebox{-20pt}{
\begin{picture}(80,40)
\fboxsep=0mm
\put(0,30){\colorbox[gray]{0.7}{\makebox(10,10){}}}
\put(10,30){\colorbox[gray]{0.7}{\makebox(10,10){}}}
\put(20,30){\colorbox[gray]{0.7}{\makebox(10,10){}}}
\put(0,40){\line(1,0){50}}
\put(0,30){\line(1,0){50}}
\put(20,20){\line(1,0){40}}
\put(30,10){\line(1,0){50}}
\put(30,0){\line(1,0){50}}
\put(0,30){\line(0,1){10}}
\put(10,30){\line(0,1){10}}
\put(20,20){\line(0,1){20}}
\put(30,0){\line(0,1){40}}
\put(40,0){\line(0,1){40}}
\put(50,0){\line(0,1){40}}
\put(60,0){\line(0,1){20}}
\put(70,0){\line(0,1){10}}
\put(80,0){\line(0,1){10}}
\end{picture}
}
\\
@VVV
\\
\raisebox{-20pt}{
\begin{picture}(80,40)
\fboxsep=0mm
\put(0,30){\colorbox[gray]{0.7}{\makebox(10,10){}}}
\put(10,30){\colorbox[gray]{0.7}{\makebox(10,10){}}}
\put(30,20){\colorbox[gray]{0.7}{\makebox(10,10){}}}
\put(0,40){\line(1,0){50}}
\put(0,30){\line(1,0){50}}
\put(20,20){\line(1,0){40}}
\put(30,10){\line(1,0){50}}
\put(30,0){\line(1,0){50}}
\put(0,30){\line(0,1){10}}
\put(10,30){\line(0,1){10}}
\put(20,20){\line(0,1){20}}
\put(30,0){\line(0,1){40}}
\put(40,0){\line(0,1){40}}
\put(50,0){\line(0,1){40}}
\put(60,0){\line(0,1){20}}
\put(70,0){\line(0,1){10}}
\put(80,0){\line(0,1){10}}
\put(20,30){\makebox(10,10){$\times$}}
\end{picture}
}
\\
@VVV
\\
\raisebox{-20pt}{
\begin{picture}(80,40)
\fboxsep=0mm
\put(0,30){\colorbox[gray]{0.7}{\makebox(10,10){}}}
\put(10,30){\colorbox[gray]{0.7}{\makebox(10,10){}}}
\put(40,10){\colorbox[gray]{0.7}{\makebox(10,10){}}}
\put(0,40){\line(1,0){50}}
\put(0,30){\line(1,0){50}}
\put(20,20){\line(1,0){40}}
\put(30,10){\line(1,0){50}}
\put(30,0){\line(1,0){50}}
\put(0,30){\line(0,1){10}}
\put(10,30){\line(0,1){10}}
\put(20,20){\line(0,1){20}}
\put(30,0){\line(0,1){40}}
\put(40,0){\line(0,1){40}}
\put(50,0){\line(0,1){40}}
\put(60,0){\line(0,1){20}}
\put(70,0){\line(0,1){10}}
\put(80,0){\line(0,1){10}}
\put(20,30){\makebox(10,10){$\times$}}
\put(30,20){\makebox(10,10){$\times$}}
\end{picture}
}
@>>>
\raisebox{-20pt}{
\begin{picture}(80,40)
\fboxsep=0mm
\put(0,30){\colorbox[gray]{0.7}{\makebox(10,10){}}}
\put(30,10){\colorbox[gray]{0.7}{\makebox(10,10){}}}
\put(40,10){\colorbox[gray]{0.7}{\makebox(10,10){}}}
\put(0,40){\line(1,0){50}}
\put(0,30){\line(1,0){50}}
\put(20,20){\line(1,0){40}}
\put(30,10){\line(1,0){50}}
\put(30,0){\line(1,0){50}}
\put(0,30){\line(0,1){10}}
\put(10,30){\line(0,1){10}}
\put(20,20){\line(0,1){20}}
\put(30,0){\line(0,1){40}}
\put(40,0){\line(0,1){40}}
\put(50,0){\line(0,1){40}}
\put(60,0){\line(0,1){20}}
\put(70,0){\line(0,1){10}}
\put(80,0){\line(0,1){10}}
\put(10,30){\makebox(10,10){$\times$}}
\end{picture}
}
\\
@VVV @VVV
\\
\raisebox{-20pt}{
\begin{picture}(80,40)
\fboxsep=0mm
\put(0,30){\colorbox[gray]{0.7}{\makebox(10,10){}}}
\put(10,30){\colorbox[gray]{0.7}{\makebox(10,10){}}}
\put(50,0){\colorbox[gray]{0.7}{\makebox(10,10){}}}
\put(0,40){\line(1,0){50}}
\put(0,30){\line(1,0){50}}
\put(20,20){\line(1,0){40}}
\put(30,10){\line(1,0){50}}
\put(30,0){\line(1,0){50}}
\put(0,30){\line(0,1){10}}
\put(10,30){\line(0,1){10}}
\put(20,20){\line(0,1){20}}
\put(30,0){\line(0,1){40}}
\put(40,0){\line(0,1){40}}
\put(50,0){\line(0,1){40}}
\put(60,0){\line(0,1){20}}
\put(70,0){\line(0,1){10}}
\put(80,0){\line(0,1){10}}
\put(20,30){\makebox(10,10){$\times$}}
\put(30,20){\makebox(10,10){$\times$}}
\put(40,10){\makebox(10,10){$\times$}}
\end{picture}
}
@>>>
\raisebox{-20pt}{
\begin{picture}(80,40)
\fboxsep=0mm
\put(0,30){\colorbox[gray]{0.7}{\makebox(10,10){}}}
\put(30,10){\colorbox[gray]{0.7}{\makebox(10,10){}}}
\put(50,0){\colorbox[gray]{0.7}{\makebox(10,10){}}}
\put(0,40){\line(1,0){50}}
\put(0,30){\line(1,0){50}}
\put(20,20){\line(1,0){40}}
\put(30,10){\line(1,0){50}}
\put(30,0){\line(1,0){50}}
\put(0,30){\line(0,1){10}}
\put(10,30){\line(0,1){10}}
\put(20,20){\line(0,1){20}}
\put(30,0){\line(0,1){40}}
\put(40,0){\line(0,1){40}}
\put(50,0){\line(0,1){40}}
\put(60,0){\line(0,1){20}}
\put(70,0){\line(0,1){10}}
\put(80,0){\line(0,1){10}}
\put(10,30){\makebox(10,10){$\times$}}
\put(40,10){\makebox(10,10){$\times$}}
\end{picture}
}
\\
@. @VVV
\\
@.
\raisebox{-20pt}{
\begin{picture}(80,40)
\fboxsep=0mm
\put(0,30){\colorbox[gray]{0.7}{\makebox(10,10){}}}
\put(40,0){\colorbox[gray]{0.7}{\makebox(10,10){}}}
\put(50,0){\colorbox[gray]{0.7}{\makebox(10,10){}}}
\put(0,40){\line(1,0){50}}
\put(0,30){\line(1,0){50}}
\put(20,20){\line(1,0){40}}
\put(30,10){\line(1,0){50}}
\put(30,0){\line(1,0){50}}
\put(0,30){\line(0,1){10}}
\put(10,30){\line(0,1){10}}
\put(20,20){\line(0,1){20}}
\put(30,0){\line(0,1){40}}
\put(40,0){\line(0,1){40}}
\put(50,0){\line(0,1){40}}
\put(60,0){\line(0,1){20}}
\put(70,0){\line(0,1){10}}
\put(80,0){\line(0,1){10}}
\put(10,30){\makebox(10,10){$\times$}}
\put(30,10){\makebox(10,10){$\times$}}
\end{picture}
}
\\
@. @VVV
\\
@.
\raisebox{-20pt}{
\begin{picture}(80,40)
\fboxsep=0mm
\put(30,0){\colorbox[gray]{0.7}{\makebox(10,10){}}}
\put(40,0){\colorbox[gray]{0.7}{\makebox(10,10){}}}
\put(50,0){\colorbox[gray]{0.7}{\makebox(10,10){}}}
\put(0,40){\line(1,0){50}}
\put(0,30){\line(1,0){50}}
\put(20,20){\line(1,0){40}}
\put(30,10){\line(1,0){50}}
\put(30,0){\line(1,0){50}}
\put(0,30){\line(0,1){10}}
\put(10,30){\line(0,1){10}}
\put(20,20){\line(0,1){20}}
\put(30,0){\line(0,1){40}}
\put(40,0){\line(0,1){40}}
\put(50,0){\line(0,1){40}}
\put(60,0){\line(0,1){20}}
\put(70,0){\line(0,1){10}}
\put(80,0){\line(0,1){10}}
\put(0,30){\makebox(10,10){$\times$}}
\end{picture}
}
\end{CD}
$$
\caption{Excited diagrams in a swivel}
\label{fig:swivel_excited}
\end{figure}
\end{exam}

Let $P$ be a connected $d$-complete poset 
and $W$ the Weyl group corresponding to the top tree $\Gamma$ viewed as a Dynkin diagram. 
Fix a linear extension of $P$, 
i.e. a labeling of the elements of $P$ with $p_1, \dots, p_N$ so that $p_i < p_j$ in $P$ 
implies $i < j$.
Then we can associate to a subset $D$ of $P$ a well-defined element $w_D \in W$ as in 
\eqref{eq:w_D}.
The following proposition gives a characterization of excited diagrams.

\begin{prop}
\label{prop:E=R}
Let $P$ be a connected $d$-complete poset and let $F$ be an order filter of $P$.
Then a subset $D \subset P$ is an excited diagram of $F$ in $P$ 
if and only if $\# D = \# F$ and $w_D = w_F$.
\end{prop}

To prove this proposition, we prepare two lemmas.

\begin{lemm}
\label{lem:comm}
Let $D$ be a subset of $P$ and $u$ and $v$ elements of $P$ such that $v < u$ and $c(u) = c(v) = i$.
Suppose $v = p_k$ and $u = p_l$.
If $[v,u] \cap D \cap N_i = \emptyset$, 
then we have $s(p_j) s(p_l) = s(p_l) s(p_j)$ for any $p_j \in D$ with $k < j < l$. 
\end{lemm}
